\section{Qué son los verbos modales}
\textbf{Función.} Modifican el significado del verbo principal: capacidad, permiso, deseo, obligación o intención.

\textbf{Lista A1.} \textit{kunnen} (poder), \textit{moeten} (deber), \textit{willen} (querer), \textit{mogen} (tener permiso), \textit{zullen} (plan/intención ligera). En este capítulo practicamos \textit{kunnen}, \textit{moeten} y \textit{willen}.

\textbf{Estructura básica}
\begin{itemize}[leftmargin=*]
  \item [S] + [Modal conjugado en V2] + \ldots{} + [Infinitivo final].
  \item Negación: \textit{niet} va antes del infinitivo. Ej.: \textit{Ik wil niet werken}.
  \item Pregunta de sí/no: el modal pasa a primera posición. Ej.: \textit{Kan jij komen?}
\end{itemize}

\textbf{Mini práctica}
\begin{itemize}[leftmargin=*]
  \item Forma una frase con modal + infinitivo: \textit{ik / kunnen / koken}.
  \item Convierte a pregunta: \textit{Jij wilt koffie}. ¿Cómo quedaría?
\end{itemize}
